\documentclass{article}
\usepackage{booktabs}
% A toy paper for the gate examples. Every number is invented.
\begin{document}

\section{Data}

The sample holds 12,480 pupils in 416 schools. Section~2 reports the estimates.

\section{Results}

Smaller classes go with higher test scores. In the baseline the coefficient on class size is $-$0.142 (Table~\ref{tab:main}, column (1)). With the instrument it is $-$0.215, significant at the one percent level. A class smaller by five pupils is associated with a score higher by 1.08 points (5 $\times$ 0.2153), about 1.8\% of the mean score.\footnote{The mean score is 61.37 points.}
% the retired run printed $-$0.207 in the sentence above

With the control variables the instrumented coefficient is $-$0.188, significant at the five percent level. In the baseline, the 95\% confidence interval of the five-pupil score gain is [0.45,\,1.70] points.

\begin{table}[h]
\centering
\caption{Class size and test scores}\label{tab:main}
% generated from results/current/estimates.csv; do not edit by hand
\begin{tabular}{lccc}
\toprule
 & (1) & (2) & (3) \\
 & OLS & IV & IV \\
\midrule
Class size & $-$0.142$^{***}$ & $-$0.215$^{***}$ & $-$0.188$^{**}$ \\
 & (0.031) & (0.064) & (0.071) \\
Controls & No & No & Yes \\
First-stage $F$ & & 38.42 & 35.10 \\
Pupils & 12,480 & 12,480 & 12,480 \\
\bottomrule
\end{tabular}

\par\footnotesize Standard errors in parentheses.\\% the earlier run reported 12,512 pupils
Stars: *** $p<0.01$, ** $p<0.05$, * $p<0.10$.
\end{table}

\end{document}
